\section{A compact contraction with Grassmannian fibre}\label{sec:construction}

We now construct a family for which the general statements have
two independent Schubert directions to detect.
Its central repeated factor has multiplicity two, so its effective
local group is \(U(2)\). The construction includes the full compact
moduli spaces, not only the Grassmannian quotient near one point.

\subsection{The surface and the complete factor list}
Choose a K3 surface with
\begin{equation}\label{eq:rankseven-lattice}
 \NS(X)=\Z H\oplus\Z D,\qquad H^2=2,\quad H\cdot D=0,\quad D^2=-44.
\end{equation}
Such a surface exists: the primitive embedding
\(H=e_1+f_1,\ D=e_2-22f_2\) into two hyperbolic planes and
Morrison's realization theorem
\cite[Corollary~2.9(i)]{morrison} give the full Néron--Severi
lattice. A root would satisfy \(h^2-22m^2=-1\), impossible modulo
eleven. We choose the positive component so that \(H\) is ample.
Put
\begin{equation}\label{eq:rankseven-vectors}
 a=(3,D,-7),\quad b=(1,0,1),\quad
 w=a+b=(4,D,-6),\quad v=2a+b=(7,2D,-13).
\end{equation}
Then
\[
 a^2=b^2=-2,\quad(a,b)=4,\quad w^2=4,\quad v^2=6,\quad
 (a,w)=2,\quad(a,v)=0.
\]

\begin{lemma}\label{lem:rank-list}
For central polystable sheaves or residual ideal bundles of total
vector \(a+(0,0,l)\), \(w+(0,0,l)\), or \(v+(0,0,l)\), with
\(l\ge0\), the only possible stable-factor vector lists are
\[
 \{a\};\qquad \{w\},\{a,b\};\qquad
 \{v\},\{a,w\},\{a,a,b\},
\]
respectively. In every case \(l=0\).
\end{lemma}

\begin{proof}
Every stable factor has vector \((r,mD,s)\), since its \(H\)-slope
is zero. Simplicity, Serre duality and Mukai Riemann--Roch give
its square at least \(-2\), hence
\[
 s\le\left\lfloor\frac{1-22m^2}{r}\right\rfloor.
\]
For total vector \((R,MD,S+l)\), summation and Cauchy--Schwarz give
\begin{equation}\label{eq:rank-bound}
 22\sum_i\frac{m_i^2}{r_i}\le
 \sum_i\frac1{r_i}-S-l\le R-S,\qquad
 \left(\sum_i|m_i|\right)^2\le\frac{R(R-S)}{22}.
\end{equation}
For \(a,w\) the last bounds are \(30/22,40/22\).
Since \(\sum m_i=1\), there is exactly one nonzero \(m_i\), equal
to one. For \(v\) the bound is \(140/22<9\), and
\(\sum m_i=2\); the nonzero entries are one two or two ones.
In particular negative entries cannot occur.

For zero-\(m\) factors of total rank \(z\), the maximal sum of
third components is \(z\), attained only by \(z\) rank-one factors
with \(s=1\). For \(a\), the maximal total third component with
the nonzero factor of rank \(r=1,2,3\) is
\(\lfloor-21/r\rfloor+3-r=-19,-10,-7\).
For \(w\), the corresponding values at \(r=1,2,3,4\) are
\(-18,-9,-6,-6\). These give precisely the stated lists.

For \(v\), a single \(m=2\) factor of rank \(r=1,\ldots,7\)
gives the bounds
\[
 -81,\ -39,\ -25,\ -19,\ -16,\ -14,\ -13.
\]
For two \(m=1\) factors, order the ranks \(r\le q\), \(r+q\le7\).
The bound is
\(\lfloor-21/r\rfloor+\lfloor-21/q\rfloor+7-r-q\).
All possibilities are
\[
\begin{array}{c|rrrrrr}
(r,q)&(1,1)&(1,2)&(1,3)&(1,4)&(1,5)&(1,6)\\ \hline
\text{bound}&-37&-28&-25&-25&-25&-25\\[2pt]
(r,q)&(2,2)&(2,3)&(2,4)&(2,5)&(3,3)&(3,4)\\ \hline
\text{bound}&-19&-16&-16&-16&-13&-13
\end{array}
\]
Only \(r=7\) in the first case, and \((3,3),(3,4)\) in the
second, survive. They give \(v,2a+b,a+w\).
Every surviving bound already equals \(S\). Thus \(l=0\) and
each factor attains its bound. This exhausts all ranks and charges.
\end{proof}

The double dual of a slope-stable torsion-free sheaf is slope-stable.
A positive discrepancy of length \(l\) lowers its Mukai square by
\(2rl\). For the vectors \(a,w,v\), the resulting squares would
be at most \(-8,-4,-8\), respectively, contradicting the stable
square bound. Thus stable sheaves of these vectors are bundles.
The vector \(b\) gives \(\OO_X\). Extensions of the listed stable
bundles are bundles, so every central semistable sheaf of these
vectors is locally free.

\subsection{Existence and compactness}
Yoshioka's generic-polarization theorem
\cite[Theorem~0.1]{yoshioka99} gives a stable object of vector
\(a\) in a chamber adjoining \(H\).
A proper positive-rank class cannot have normalized first Chern
class \(D/3\). After avoiding the finite slope walls, generic
Gieseker stability is therefore slope stability.
Move the polarization to \(H\) inside this chamber. Each fixed
subsheaf inequality becomes weak, so the object is \(H\)-semistable.
Lemma~\ref{lem:rank-list} makes it stable, and the double-dual
argument makes it locally free. Denote it by \(A\).
If two nonisomorphic \(H\)-stable bundles had vector \(a\), their
\(\Hom\) and \(\Ext^2\) would vanish, contradicting Euler
characteristic \(-(a,a)=2\). The same identity gives
\(\Ext^1(A,A)=0\). Thus \(A\) is unique and spherical.
Set \(B=\OO_X\).

The primitive vectors \(w,v\) similarly have compact generic
moduli of dimensions six and eight, of deformation types
\(X^{[3]}\) and \(X^{[4]}\)
\cite[Appendix, Theorem~8.1]{yoshioka01}.
For \(v\), use large integral polarizations \(mH\pm D\):
the degree \(\mp88\) is coprime to seven.
For \(w\), avoid the locally finite slope walls; equality of
normalized first Chern classes would require \(rD/4\) integral
with \(0<r<4\), which is impossible.
The double-dual argument applies at any polarization and makes
every slope-stable object in these compact moduli locally free.
Write \(S_\pm\) for the sixfolds and \(M_\pm\) for the eightfolds.

Finite wall theory \cite[Theorem~6.6]{grebtoma} permits a sufficiently
small interval \(H+tD\) with no wall except possibly \(t=0\).
Every chamber object is \(H\)-semistable by the preceding limiting
inequalities. Conversely every \(H\)-stable object persists in both
adjacent chambers by openness and chamber constancy.
The local models below show that \(t=0\) is an actual wall.
Let \(\omega_t\) be the smoothly varying Ricci-flat representatives
of these classes.

The eightfold is fine: the algebraic \(K\)-class
\(\OO_X+\OO_p\) has index \(\chi(E)+\rank E=-6+7=1\);
the universal-family criterion applies
\cite[Theorem~4.6.5]{huybrechtslehn}.
For \(w\), algebraic Mukai divisibility is
\(\gcd(4,44,6)=2\). Since \(w^2=4>0\), the positive-square
obstruction theorem of Mattei--Meinsma \cite{matteimeinsma}
gives Brauer order two. We use local or twisted families on
\(S_\pm\), without assuming an analytic trivialization of this
obstruction. Topological universal classes remain available for
the index calculations.

Finally Lemma~\ref{lem:rank-list} excludes positive-charge central
ideal objects. Surface compactness and removability
\cite[Section~3.1]{gstw}, with the uniform small-energy estimates
for the smooth family \(\omega_t\), give smooth subsequential
gauge convergence on all of \(X\) as \(t\to0\).
There are no positive-charge ideal objects at sufficiently nearby
parameters either. Otherwise choose such parameters tending to
zero. A bounded-energy subsequence has a fixed positive defect
integer; compactness of the residual connections adds only
nonnegative defect, contradicting the central list.
This verifies the full no-boundary requirement, including nearby
parameters, rather than just local freeness on the displayed fibres.

\subsection{The full nonabelian local model}\label{subsec:grass-model}
At \(E_o=A\otimes\C^2\oplus B\), both self-\(\Ext^1\) spaces
vanish and the cross dimension is four.
Let \(\mathbb V=\Ext^1(B,A)^*\simeq\C^4\), and use coordinates
\[
 Z_1:\mathbb V\to\C^2,\qquad Y_1:\C^2\to\mathbb V.
\]
The subscript distinguishes this linear map from the global
target \(Z\). The effective group is \(U(2)\), and the equations are
\begin{equation}\label{eq:u2-model}
 Z_1Y_1=0,\qquad Z_1Z_1^*-Y_1^*Y_1=\zeta I_2,\qquad
 \zeta=\frac{88\pi}{7}t .
\end{equation}
Indeed \eqref{eq:levels} gives
\(2\pi(3\mu_t(v)-\deg_t A)=88\pi t/7\);
the \(B\)-vertex level is \(-2\zeta\).
These are full convergent coordinates by the normal form and
harmonic reduction stated in Section~\ref{sec:data}.
Qualitative formality is already covered by
Bandiera--Manetti--Meazzini \cite[Theorem~5.1]{bmm}: the
automorphism group is reductive, although the stable factors
have different reduced Hilbert constants.
The additional coordinate estimates, rather than formality
alone, supply the relative harmonic comparison.

For \(\zeta>0\), \(Z_1\) is surjective. Set \(P=\ker Z_1\),
\(\mathcal U_P=\mathbb V/P\), and \(N=Y_1Z_1\).
Then \(\im N\subset P\subset\ker N\).
Conversely such an incident pair determines \(Z_1,Y_1\)
up to \(GL(2)\). Thus the quotient is
\begin{equation}\label{eq:incidence}
 T^*\Gr(2,\mathbb V)=
 \{(N,P):\im N\subset P\subset\ker N,\ \dim P=2\}
 \longrightarrow
 \mathcal N=\{N:N^2=0,\ \rank N\le2\}.
\end{equation}
The cotangent zero section \(G=\Gr(2,\mathbb V)\) has
\(Y_1=0\), \(Z_1Z_1^*=\zeta I_2\), and coefficient norm
squared \(2\zeta\).
For \(\zeta<0\), \(Y_1\) is injective and \(P=\im Y_1\)
gives the same incidence resolution over \(\mathcal N\).
Thus this equal-dimensional Grassmannian model is not a
nontrivial flop. Fibres over ranks two, one and zero are,
respectively, a point, \(\PP^1\) and \(G\).

For clarity we check normality of the affine target used here.
The variety \(\{Z_1Y_1=0\}\subset\C^{16}\), over
\(\rank Z_1=k\), has dimension
\(k(6-k)+2(4-k)=8+4k-k^2\), with maximum twelve at \(k=2\).
Every lower-rank pair is a limit of rank-two pairs by perturbing
\(Z_1\), with \(Y_1\) fixed, among maps annihilating \(\im Y_1\).
It is an irreducible complete intersection of four equations.
If either map has rank two, the differential of \(Z_1Y_1\)
is onto \(\End(\C^2)\). The locus where both ranks are at most
one has dimension at most nine. The complete intersection is
therefore regular in codimension one and normal.
Reductive invariants preserve normality, and the matrix invariant
\(Y_1Z_1\) identifies the quotient with \(\mathcal N\).

The model at \(A\oplus B\), for vector \(w\), is the four-arrow
\(U(1)\) quotient with level \(22\pi t\). It contracts
\(T^*\PP^3\) to the rank-at-most-one square-zero cone.
At \(A\oplus F\), for an \(H\)-stable \(w\)-bundle \(F\),
the self-deformation space of \(F\) has dimension six and
\(\dim\Ext^1(F,A)=(a,w)=2\).
The local contraction is consequently a smooth sixfold factor
times \(T^*\PP^1\to\C^2/\{\pm1\}\), with transverse level
\(\zeta\).

The simple-zero condition is explicit. In \eqref{eq:u2-model}
at zero level, take
\(Z_1=(I_2\ 0)\), \(Y_1=(0\ I_2)^{\mathsf T}\).
A subrepresentation with zero \(B\)-component has its
\(A\)-component annihilated by all rows of \(Y_1\), hence is zero.
A nonzero \(B\)-component forces the full \(A\)-component by the
columns of \(Z_1\). The auxiliary model uses the analogous pair
\((1,0,0,0)\), \((0,1,0,0)^{\mathsf T}\).
The two-arrow transverse model has the same two distinct
nonzero coordinates. Trivial tangent directions do not change
simplicity. Small multiples in the convergent charts give actual
\(H\)-stable bundles. Their loci are nonempty and dense in the
irreducible chamber moduli.

\subsection{Whole fibres and normal targets}\label{subsec:targets}
The slope classifying theorem of Greb--Toma
\cite[Main Theorem, Theorem~5.5 and Remark~5.9]{grebtoma}
first gives maps to the common reduced slope images.
For the auxiliary family the local classifying maps descend
despite base-line changes. At this stage we do not assume
normality of the image or its analytic identification with
\(\mathcal N\).

We verify whole-fibre coverage before making that identification.
Suppose \(t_n\to0\) and a chamber bundle \(E_n\) has central
graded object \(2A\oplus B\). The no-charge result gives a smooth
gauge limit \(E\). In a Jordan--Hölder filtration choose its
first inclusion and last quotient, pass to fixed factor types,
and normalize these holomorphic maps separately to \(L^2\)-norm
one in the converging gauges. Elliptic estimates give nonzero
limiting maps. A nonzero map between stable equal-slope factors
splits the corresponding factor in the polystable bundle \(E\).

If the first and last factors differ, their limits split off
\(A\) and \(B\); the remaining vector \(a\) gives the unique \(A\).
If both are \(A\), the inclusion \(i_n\) and quotient \(q_n\)
satisfy \(q_ni_n=0\), still true after separate normalization
and passage to the limit. With only one \(A\)-summand in \(E\),
both nonzero maps would factor through it and have nonzero
composition. Thus there are at least two copies of \(A\);
the remaining vector is \(b\). In all cases \(E=2A\oplus B\).
The two-factor proof gives whole-fibre coverage for \(A\oplus B\)
and for \(A\oplus F\) with \(F\) in a compact part of its stable
locus.

We also need joint coverage when \(F_n\to A\oplus B\).
At fixed \(H\), the surface comparison
\cite[Theorem~4.11 and Corollary~4.15]{gstw} identifies the
slope and Uhlenbeck topologies. Since charge is zero, \(F_n\)
converges in gauge to \(A\oplus B\).
Orient the extensions of \(A\) and \(F_n\) in either chamber and
normalize their two maps as above. Their limiting composition
vanishes. If the limiting map involving \(A\oplus B\) has a
nonzero \(B\)-component, it splits off \(B\) and \(A\).
Otherwise zero composition forces two \(A\)-summands.
The limit is again \(2A\oplus B\).
One fixed deep chart therefore covers these joint limits.

All actual fibres lie in these charts for sufficiently small
parameters. Conversely every finite core has norm \(O(\sqrt{|t|})\)
and lies in the chart. They classify the same holomorphic bundles
and central graded classes, so chamber constancy identifies
their whole fibres.

Normalize the common reduced projective image. Connectedness of
the fibres just identified and Stein factorization make this
normalization bijective on points. Denote the resulting targets
by \(Z\) for \(v\) and \(T\) for \(w\).
Whole-fibre coverage and properness give saturated neighborhoods
in the actual and finite resolutions. Pull back target holomorphic
functions through the classifying charts and descend them using
\(\pi_*\OO=\OO\) for a proper birational map to a normal target.
Doing this in both directions gives inverse germ maps: their
composites agree on the dense stable locus and hence everywhere.
Thus the contractions have exactly the analytic models above.

The target \(T\) has one singular point \(o\), represented by
\(A\oplus B\); \(S^\circ=T\setminus\{o\}\) is the stable
\(w\)-moduli. The map \(F\mapsto A\oplus F\) embeds it into \(Z\).
At \(o\), fix a spectator \(A\)-copy in the \(U(2)\) model.
The equivariant Kuranishi and Darboux maps preserve this fixed
locus, which is the \(U(1)\) four-arrow representation.
On invariants its inclusion imposes the two-by-two minors of
\(N\), giving the closed rank-at-most-one subcone of \(\mathcal N\).
These local embeddings glue because they classify the same
direct sums. Properness makes \(T\to Z\) a closed algebraic
embedding. The singular support of \(\mathcal N\) is its
rank-at-most-one locus. Consequently \(T\) is precisely the
\emph{reduced} singular locus of \(Z\).

We have proved projective symplectic contractions
\(\pi_\pm:M_\pm\to Z\) with fibres
\[
 \begin{array}{c|ccc}
 \text{central class}&H\text{-stable }v&A\oplus F,\ F\in S^\circ&2A\oplus B\\ \hline
 \text{fibre}&\text{point}&\PP^1&\Gr(2,4).
 \end{array}
\]
The auxiliary contractions \(S_\pm\to T\) have fibre \(\PP^3\)
over \(o\). The symplectic forms agree with the quotient forms
in the charts and with the common stable form elsewhere.
Surface comparison gives the projective/gauge topology, while
the earlier factor and charge argument excludes every further
boundary. This verifies (G1)--(G3) for the eightfold.

\subsection{The actual extension and specialization}\label{subsec:specialization}
Let \(\mathcal U\) be the positive tautological quotient on
\(G=\Gr(2,\mathbb V)\). In the positive chamber, the actual
universal family restricts, up to a base line, to
\begin{equation}\label{eq:actual-extension}
 0\longrightarrow A\boxtimes\mathcal U
 \longrightarrow\mathcal E_G
 \longrightarrow B\boxtimes1\longrightarrow0 .
\end{equation}
To see that this is the actual extension family, not just its
linear approximation, observe that on \(Y_1=0\) the central
\(\C^*\subset GL(2)\) gives \(Z_1\) weight \(1\) and \(Y_1\)
weight \(-1\). An equivariant holomorphic coordinate change
tangent to the identity has no nonlinear term surviving on
this axis in either output. It fixes the axis. The harmonic
recursion is also linear there, since two one-direction
extension terms have zero product. The quotient
\(\mathbb V\to\mathcal U\) therefore gives
\eqref{eq:actual-extension} itself.

Take \(N_0\ne0\) with \(N_0^2=0\), \(\rank N_0=1\).
Put \(L_0=\im N_0\), \(K_0=\ker N_0\).
The entire fibre over \(\tau N_0\), \(\tau\ne0\), is
\[
 \{P:L_0\subset P\subset K_0,\ \dim P=2\}\simeq\PP^1.
\]
Its closure is the product family
\[
 \PP^1\times\Delta\longrightarrow T^*G,\qquad
 (P,\tau)\longmapsto(\tau N_0,P).
\]
It is proper and flat, with reduced central fibre a minimal
Schubert line in \(G\). The whole-fibre chart transfers it to
the compact moduli. Thus the moving fibre class \(f\) specializes
with multiplicity one to the line class \(\ell\):
\begin{equation}\label{eq:line-specialization}
 f\longmapsto\ell .
\end{equation}
\begin{figure}[tb]
 \centering
 \begin{tikzpicture}[>=Latex, every node/.style={font=\small},
                    box/.style={align=center,inner sep=5pt}]
 \node[box] (g) at (0,2.8) {\(G=\Gr(2,\mathbb V)\)};
 \node[box] (l) at (-2.6,1.4)
   {\(\Sigma_L=\{P:L_0\subset P\}\)\\ \(\simeq\PP^2\)};
 \node[box] (k) at (2.6,1.4)
   {\(\Sigma_K=\{P:P\subset K_0\}\)\\ \(\simeq\PP^2\)};
 \node[box] (line) at (0,-0.15)
   {\(\{P:L_0\subset P\subset K_0\}\simeq\PP^1\)};
 \draw[->] (l) -- (g);
 \draw[->] (k) -- (g);
 \draw[->] (line) -- (l);
 \draw[->] (line) -- (k);
 \node[box] at (0,-1.05)
   {\(\dim L_0=1,\quad\dim K_0=3,\quad L_0\subset K_0\subset\mathbb V\)};
\end{tikzpicture}

 \caption{The Schubert line obtained from \(N=\tau N_0\).
 For the fixed incident flag \(L_0\subset K_0\), the two
 Schubert planes meet along this line. The chosen flat family
 has this line as its reduced central fibre inside the larger
 fibre \(G\). Arrows denote inclusions.}
 \label{fig:schubert}
\end{figure}
The directions are parametrized by flags
\(L_0\subset K_0\subset\mathbb V\); including the plane gives
\(L_0\subset P\subset K_0\).
Over \(S^\circ\), projective extension transitions are in
\(PGL_2\), so their action on the fibre's integral cohomology
is trivial. This does not trivialize the projective bundle.

The normal bundle of \(G\) is \(T^*G\).
Writing a cotangent vector as a map
\(\mathbb V/P\to P\), the preimage of the rank-at-most-one locus has
equation given by its \(2\times2\) determinant.
Its restriction line bundle is
\(\det P\otimes\det(\mathbb V/P)^*=\OO_G(-2h)\), where
\(h=c_1(\mathcal U)\). This records the incidence of the moving
exceptional divisor with the deep core.
